\paragraph{生産者物価水準目標（PPLT）}
\label{sec:chapter5_a_shock_policies_pplt}
PPLTの式については
\ref{sec:chapter5_beta_shock_policies_pplt}
を参照のこと。
\begin{enumerate}
    \item
    \label{itm:chapter5_a_shock_policies_pplt_c_H_W}
        \(c_t^{H\to W}\)：
        総消費指数\(c_t^{H\to W}\)の図\ref{fig:chapter5_a_shock_graphs_c_H_W}を見ると、PPLTの軌道は消費者物価水準目標（CPLT）や消費者物価インフレ目標（IT-CPI）などと同様に、ショック直後の初期において極端な大暴落を引き起こしていることがわかる。
        総消費水準目標（CLT）や名目総消費水準目標（NCLT）が落ち込みをほぼゼロに抑え込んでいるのとは対照的である。
        生産者物価水準を一定に保とうとする強力な制約が実体経済への強い下押し圧力として働き、消費の深刻な減少を招いている。
        消費の低下は厚生にマイナスに働くため、この前半における深刻な消費の落ち込みがPPLTの厚生を全体の中で下位クラスへと低迷させる大きな要因となっている。
        また深く落ち込んだ状態から極めて長い時間をかけて定常状態へ回帰するだけの軌道となっているため、消費増加による厚生の改善効果も得られていない。
    \item
    \label{itm:chapter5_a_shock_policies_pplt_y_H}
        \(y_t^H\)：
        生産\(y_t^H\)の図\ref{fig:chapter5_a_shock_graphs_y_H}を見ると、PPLTの軌道は総消費指数\(c_t^{H\to W}\)と同様に、消費者物価水準目標（CPLT）や生産者物価インフレ目標（IT-PPI）と極めて類似した極端な大暴落を見せていることがわかる。
        生産者物価水準を一定に保とうとする強力な制約が働くため、供給ショックに対して生産を極端に犠牲にすることで対応しており、ショック直後の前半においてマイナス方向へ極めて大きく落ち込んでいる。
        生産の低下は労働の減少を意味するため、この前半部分においては労働の非効用を減らして厚生にプラスに働く効果が他の政策よりも大きくなっている。
        しかしながら、前述した深刻な消費の大暴落によるマイナス効果を補うほどの改善には至っておらず、結果として全体的な厚生の評価を低迷させる結果となっている。
        また深く落ち込んだ状態から極めて長い時間をかけて定常状態へ回帰するだけの軌道を描いている。
    \item
    \label{itm:chapter5_a_shock_policies_pplt_pi_H}
        \(\pi_t^H\)：
        グロス・インフレ率\(\pi_t^H\)の図\ref{fig:chapter5_a_shock_graphs_pi_H}を見ると、PPLTの軌道は生産者物価インフレ目標（IT-PPI）と同様に、全期間を通じてゼロ近傍で極めて平坦な推移を完全に維持していることがわかる。
        総消費水準目標（CLT）や名目総消費水準目標（NCLT）などが初期に大きなプラス方向へのスパイクを記録しているのや、消費者物価水準目標（CPLT）などが極端なマイナス方向（デフレ）への下落を引き起こしているのとは対照的である。
        生産者物価水準を目標経路に厳格に固定しようとする強力なメカニズムが働くため、供給ショックに伴う物価変動圧力を完全に封じ込め、国内生産財のインフレ率の変動を完全に抑制している。このインフレ率の完全な安定化が、次項で述べる価格分散の蓄積を回避する最大の要因となっている。
    \item
    \label{itm:chapter5_a_shock_policies_pplt_Delta_t_minus_1_H}
        \(\Delta_{t-1}^H\)：
        価格分散\(\Delta_{t-1}^H\)の図\ref{fig:chapter5_a_shock_graphs_Delta_H}を見ると、PPLTの定常状態からの乖離は、生産者物価インフレ目標（IT-PPI）と同様に全期間を通じてほぼゼロの水準を完全に維持していることがわかる。
        総消費水準目標（CLT）や生産水準目標（OLT）などが極めて大きく増大しているのとは対照的である。
        これは前述のインフレ率の完全な安定化により、価格改定に伴う相対価格のばらつきが一切生じなかった結果であり、この面においては生産効率の悪化による厚生へのマイナス効果を完全に防いでいる。
        しかしながら、供給ショック（\(a\)ショック）においては価格分散の増大を抑制することのプラス効果以上に、実体経済（消費や生産）の大暴落によるマイナス効果が極めて絶大である。そのため、価格分散の発生を完全に抑え込んだにもかかわらず、初期の深刻な不況を防げなかったことが致命傷となり、PPLTの厚生は全体の中で下位クラスへと低迷する結果となっている。
\end{enumerate}